\section{Improving a road while preserving locality}
\label{improvement:section}

The link estimates supply inexpensive crossings after clipping, whereas
an admissible road needs a bound on its full expected cost. We bridge this
gap in two stages. First, compare each crossing with the corresponding
piece of the input road. Then replace clusters of large positive costs by
pieces of an independent spare road. The replacement tests use a fixed
number of neighboring slots. After a fixed enlargement of the spatial
scale, the resulting path belongs to the original road class.

We construct the improved $A$-road; interchanging the two families gives
the $B$-road. Positions are in $\mathcal G_n$ and costs remain in their
original units. Write $(x,y)$ for coordinates in the basis $(v,B)$;
the horizontal coordinate $x$ was denoted by $u$ in
Section~\ref{links:section}. Put
\[
 a_0=d=1-s^2,\qquad A=(a_0,s),\qquad B=(0,1)
\]
and take $0<s<1/10$. The tracking and locality norm is still the maximum
norm in $(v,w)$ coordinates. The change
$(x_v,x_w)\mapsto(x_v-sx_w,x_w)$ and its inverse have norm at most
$1+s\le1.1$; this factor is included in every band clearance below.
Scaled lattice edges have length at most $2K_{\rm lat}/n$ in these
coordinates, so their rounding errors fit the stated strict clearances
once $n$ is large. Constants are uniform over the fixed neighborhoods of
the flat endpoints. Euclidean metric estimates may use constants
depending on the hypothetical flat.

\subsection{Two inputs to the construction}
\label{improvement:interfaces}

Fix an input pair whose mean period excesses are bounded by $\beta$, where
\[
 c_0\leq\beta\leq C_s n^{1-c}.
\]
We use the preceding results in the following two forms. The gauge
input (G) combines Lemmas~\ref{found:lem-finite-cell},
\ref{found:lem-envelope}, \ref{found:lem-common-gauge}, and
\ref{phases:locality}, with the interpolation identity
\eqref{found:eq-short}. The link input (L) combines
Propositions~\ref{links:asymmetric} and \ref{links:matched} with
\eqref{links:negative-moment}. Stating their support bounds here will make
the eventual rescaling independent of all auxiliary proof parameters.

\begin{enumerate}
\item[(G)] In an iid joint world, all specified road anchors have
integrable local gauges, with the deterministic means of their label
types included. Corrected increments along a bounded number of road
periods, and changes of labels at the specified intersections, have
$L^1$ bounds $C\beta$. Comparisons with increments of the calibration
have the same bound. These comparisons hold for the supremum over
attachment occurrences in the indicated bounded parameter intervals.
Along each unit road period the gauge is interpolated so that its
corrected signed cost over a nominal interval of length $u$ is $u$
times that period's corrected increment. Consequently, over an
interval of length at most $c_1s$, situated in a prescribed bounded
number of periods, its absolute corrected cost is bounded by
$c_1s$ times a local variable of mean $C\beta$.

A baseline generated from its tube and private completion admits a
compatible gauge depending only on these common inputs. In each joint
world containing that baseline the difference from its world gauge,
including the supremum over bounded attachment intervals, has
$L^1$ norm at most $C\beta$. Every computed gauge, corrected road increment, and gauge-to-gauge
seam record has a fixed support radius $R_g$, independent of
$s,n,\beta$, probability tolerances, phase-group size, and the chosen
input road rule. Comparisons involving the calibration itself are
analytic $L^1$ bounds; those calibration errors are not asserted to
be local.

\item[(L)] Between the walls $x=ja_0$ and $x=(j+1)a_0$, let $X$ be
the minimum signed link cost, corrected at the two wall center levels,
in the horizontal band of half-width $s/50$. The admissible crossings
have horizontal drawdown at most $s/20$, and the signed endpoint arms
are along the labeled walls. The link is a function of the inputs in
a fixed neighborhood of this band. Its negative part satisfies
$\mathbb E(X/\beta)^-\leq C_X$. For every fixed upper cap $K$ and
$\eta>0$, its asymmetric clipped mean is at most $\eta$ after a
sufficiently large fixed lower cap is chosen. All direction tolerances and lower thresholds for $n$ are uniform
over the input road rules satisfying the displayed common bounds,
and over the finitely many row and phase versions used below.

The matched-cap conclusion is also available. Given $M_0<\infty$,
$\eta>0$, and finitely many additional nonnegative local variables
with uniformly bounded first moments, there is a fixed finite dyadic
set $\mathcal M\subset[M_0,\infty)$ such that at each sufficiently
large admissible $n$ a deterministic $M\in\mathcal M$ satisfies
\begin{equation}
 \mathbb E\operatorname{clip}_M(X/\beta)\leq\eta,
 \qquad
 M\sum_\ell\mathbb P(U_\ell>M)\leq\eta.
 \label{improvement:interfaces:eq1}
\end{equation}
Here $\operatorname{clip}_M(z)=f_{M,M}(z)=\max\{-M,\min\{z,M\}\}$.
The geometry and threshold for $n$ are chosen uniformly over this
finite set. This is Proposition~\ref{links:matched},
not a consequence merely of asymmetric clipping.
\end{enumerate}

The support assertions in (G) and (L) refer to all the weights and
private random inputs used to calculate a variable, including the
inputs to its gauges. Finite averaging of gauge functions over mesh
offsets preserves both the estimates and a fixed support radius:
all offsets lie in one bounded fundamental cell, irrespective of the
number of offsets.

\subsection{Independent baselines and completed environments}
\label{improvement:worlds}

Only the weights on a path's traveled edges need to be the actual
weights. We may therefore complete the environment privately outside a
deterministic band, while retaining the correct iid law for every local
calculation. We arrange the completions so that the entire primary field
is independent of the spare baseline, and so that transitions can use
both baselines with compatible gauges.

Put $z(x,y)=y-sx/a_0$. We use the deterministic primary band
\[
 \mathcal T_p=\{|z|\leq3s\}
\]
and the narrow spare-baseline tube around $z=10s$. The input tracking
bound $r=s/1000$ places the primary and spare baselines in disjoint
subsets of these bands. Edges meeting their boundaries can be assigned
by a fixed convention; for sufficiently large $n$ the $O(1/n)$ lattice
padding leaves the two preserved sets disjoint.

Generate each baseline once from the actual weights in its narrow tube,
its own independent private exterior, and its private road seeds. Its
physical path stays in that tube. Whenever a completed environment uses
this baseline, insert that same path with those same private inputs.
Choose these baseline private arrays jointly independent of the actual
edge field and of all world completion and auxiliary arrays. Denote the
compatible baseline gauges from (G) by $H_p$ and $H_s$.

The primary world uses actual weights in $\mathcal T_p$ and fresh iid
weights outside $\mathcal T_p$. Its auxiliary road seeds and exterior
completion are independent of all spare-baseline inputs. The primary
baseline is inserted using its already fixed private completion.
Although that completion is not the primary world's exterior, it is
an admissible independent randomization for a path whose physical
edges all have the primary world's weights. Keep the canonical wall
network and wall gauges of (L) unchanged, and attach the common
baseline as an extra path label using (G). In particular its gauge
does not replace any gauge defining $X$. Do the same with both
common baselines in every transition world. The link variables
therefore retain exactly their laws in (L), while (G) supplies the
new endpoint-label comparisons. The spare world is
constructed in the same way around the spare tube. Thus the entire
primary field constructed below is independent of the entire spare
baseline and its compatible gauge.

For each prospective transition make a joint world using actual
weights in
\[
 \mathcal T_t=\{-3s\leq z\leq13s\}
\]
and a fresh iid exterior. Insert both common baselines, with their
original private completions. Use fresh auxiliary road seeds. Distinct
transition worlds may share the actual weights in their preserved
bands; their exterior completions are independent. We never equate
calibrations from two different worlds. Comparisons at their seams
use (G).

Apply (G) separately to each baseline with only its own narrow tube
prescribed: take $\eta_0$ to be that tube's actual weights and take
$\eta_{\rm other}$ empty. In each world $j$, let $\zeta^j$ be the entire
weight array on the complement of this tube, using the actual weights
where that world retains them and its fresh exterior weights elsewhere.
The partition is deterministic, so each hybrid complement has the iid
exponential law and is independent of $\eta_0$. The baseline's private
field is jointly independent of $(\eta_0,\zeta^j)$. The complements in
different worlds may share actual coordinates. Applying the lemma first
to the primary baseline in its primary and transition worlds, and then
to the spare baseline in its spare and transition worlds, gives the
common gauges $H_p,H_s$ used above.

\begin{table}[htbp]
\centering
\small
\begin{tabular}{@{}p{0.9in}p{1.7in}p{1.4in}p{1.4in}@{}}
\toprule
World & Actual weights retained & Inserted baselines & Compatible gauges \\
\midrule
Primary & $\mathcal T_p$ & primary & $H_p$ \\
Spare & narrow spare tube & spare & $H_s$ \\
Transition & $\mathcal T_t$ & primary and spare & $H_p,H_s$ \\
\bottomrule
\end{tabular}
\caption{The completed environments and the inputs they share. Each
world uses fresh iid weights outside its retained region. Every inserted
baseline keeps its fixed private exterior and road seeds. A compatible
gauge is compared with that baseline's gauge in the indicated world;
wall-to-baseline switches are handled within that world.}
\label{improvement:table-worlds}
\end{table}

The following phase convention makes the slot construction stationary. If $\mathcal H$ is a finite subgroup
of $\mathbb R/\mathbb Z$ containing $s$ and all the row-mesh increments
used in the construction, let $\phi_j$ be iid uniform in $\mathcal H$.
Give wall $j$ the relative phase $\phi_j$ in the rule of
Section~\ref{phases:section}; its absolute marked-height residue is
\begin{equation}
 \theta_j=\phi_j+js\pmod1.
 \label{improvement:worlds:eq1}
\end{equation}
The translation by $A$ takes wall $j$ to wall $j+1$ and adds $s$ to its
vertical coordinate. Consequently ordinary translation of the iid
sequence $(\phi_j)$ gives exactly the required translation of the
wall family, including its marked cuts. A uniform wall phase is
needed only on the bounded wall portions queried by this construction;
its seed is stored at the corresponding primary slot. Different
transition worlds receive their own such seeds. No random phase is
shared by all slots. Gauge functions, rather than a single random
mesh choice, are averaged over mesh offsets. The number of phases
can depend on the previously fixed rational mesh and $s$; their
physical storage locations and the union of their support sets remain
in a fixed neighborhood.

For the copy with nominal origin zero, group the globally indexed
primitives of Section~\ref{phases:section} into slabs indexed by
$j=\lfloor x/a_0\rfloor$, resolving boundary ties deterministically.
The coordinates retain their structural roles and road or world origins;
translated copies carry translated origins and slabs. Each slab contains
countably many independent auxiliary coordinates, with separate
coordinates for the baselines, the primary completion, and each
translated transition type, and with the same marginal law in every
translated slab for each role. The global indexing supplies the
measure-preserving action of every lattice shift; the slab grouping
here establishes locality and separation. All quantities used for slot $i$
are measurable with respect to the slab inputs with $|j-i|\leq R_*$,
for a fixed integer $R_*$. Only bounded transverse subsets of these
slabs are queried. The retained actual primary and spare edges are
disjoint, and their auxiliary coordinates are disjoint. These facts
establish unconditional independence of the two entire fields, and
independence of a slot variable from a remote transition test. They
are stronger than conditional finite dependence given phases.

\subsection{The primary choice in each slot}
\label{improvement:primary}

A slot has two physical candidates with the same endpoints: the baseline
piece and a horizontal crossing reached by forward wall travel. Their
common endpoint convention lets us choose the cheaper candidate before
estimating its cost.

Every minimizing link used in this section is the vertex-ended
representative constructed after \eqref{links:definition}.

Choose the primary baseline's common-vertex intersection $c_i$ with
wall $i$ by the local occurrence rule of Lemma~\ref{found:lem-network}.
All such intersections have
baseline parameter $i+O(r)$ and wall height $is+O(r)$, by the tracking
bounds. The permissible parameter intervals for different walls are
disjoint and ordered. The choice is therefore local and commutes with
translation by $A$. Ties between multiple occurrences are resolved
in their order along the baseline. Let $g_i$ be the primary world's
wall gauge at $c_i$.

The two competing paths from $c_i$ to $c_{i+1}$ are the baseline
subpath and the midpoint link, whose row is $(i+1/2)s$, with its
attachments along the two walls. Every left attachment is forward
because
\[
 is+O(r)<(i+1/2)s-s/50-O(r),
\]
and every right attachment is forward for the analogous reason.
The link is followed in its crossing direction. Thus both candidates
are physical paths with positive edge traversal costs. The signed
arms in the definition of $X_i$ are only bookkeeping: they cancel
inside these forward attachments.

Let $B_i$ and $L_i$ be the normalized corrected costs of these two
paths, both in the endpoint convention $g_i,g_{i+1}$. By the linear
interpolation in (G),
\begin{equation}
 L_i=X_i/\beta+E_i,\qquad \mathbb E|E_i|\leq C_1s.
 \label{improvement:primary:eq1}
\end{equation}
Indeed the algebraic difference consists exactly of the two signed
wall movements from the baseline intersections to the midpoint
center levels. Each has nominal length $O(s)$ in a bounded number of
wall periods. This estimate precedes, and does not condition on, any
later path selection.

Choose the cheaper path and put $Z_i=\min(B_i,L_i)$. This is one
common endpoint convention, with wall gauge $g_i$ shared by its two
neighboring slots. It is essential to retain these gauges here:
replacing them slot by slot by $H_p$ would add errors of order
$\beta$, whereas \eqref{improvement:primary:eq1} is of order $s\beta$.
The compatible baseline gauge is used only at the rare seams below.

Both tails of $Z_i$ have a uniform first-moment bound. For the negative
tail, calibration domination leaves only the gauge error between $c_i$
and $c_{i+1}$; the attachment supremum in (G) bounds its expectation by
$C\beta$ before the candidate is selected. For the positive tail, the
baseline competitor has defect bounded by a fixed number of complete
period defects, because its endpoints lie in fixed parameter windows.
The endpoint gauge comparison costs a further $C\beta$. After normalization,
\begin{equation}
 \mathbb E|Z_i|\leq C_0.
 \label{improvement:primary:eq2}
\end{equation}
The complete-period estimate applies even though the endpoints were
chosen from the baseline and wall environment.

Since clipping is monotone and $1$-Lipschitz,
\begin{equation}
 \mathbb E\operatorname{clip}_M(Z_i)
 \leq\mathbb E\operatorname{clip}_M(X_i/\beta)+C_1s.
 \label{improvement:primary:eq3}
\end{equation}
Apply the matched-cap estimate with the additional variable $Z_i^+$.
For any prescribed $M_0$ and $\eta$, a cap in a fixed finite dyadic
range therefore gives
\begin{equation}
 \mathbb E\operatorname{clip}_M(Z_i)\leq C_1s+\eta,
 \qquad M\mathbb P(Z_i>M)\leq\eta.
 \label{improvement:primary:eq4}
\end{equation}
The cap depends on the input law and the scale, but is deterministic
before the environment is sampled. This dependence adds no spatial
input to the rule.

All physical pieces just constructed lie in $|z|\leq s$, after
increasing the lower threshold for $n$. Thus their actual costs in
the primary world are their actual costs in the original environment.
Their full calculation, including the wall gauges, is local in the
sense established in Section~\ref{improvement:worlds}.

\subsection{Forward transitions to the spare road}
\label{improvement:transitions}

A replacement must leave the primary baseline and return to it using
physical forward travel. We reserve a fixed number of slots for each
transition. Its geometry is fixed; only the numerical success threshold
will depend on the desired failure probability.

Fix the transition length $T_*=64$. For each
integer $k$, construct a candidate from $c_k$ to the marked spare
cut at nominal parameter $k+T_*$. Follow the primary baseline to its
common-vertex intersection with the joint world's wall $k+1$, chosen
by the same local rule, cross successive
walls through midpoint bands, and after wall $k+T_*-1$ follow the
spare baseline forward to that marked spare cut. For the intervening
link between walls $j$ and $j+1$, use row
\begin{equation}
 m_j=(j+1/2)s+10s\frac{j+1/2-k}{T_*},
 \qquad k+1\leq j\leq k+T_*-2.
 \label{improvement:transitions:eq1}
\end{equation}
For the transition from the spare cut at $k$ to $c_{k+T_*}$, use instead
\begin{equation}
 m_j=(j+1/2)s+10s\left(1-\frac{j+1/2-k}{T_*}\right).
 \label{improvement:transitions:eq2}
\end{equation}
All the rows belong to a fixed finite rational mesh. At interior walls
the row increment is $s(1\pm10/T_*)$. At the first and last attachment,
the smaller available clearance is
\[
 \left(\frac12-\frac{15}{T_*}\right)s
 =\frac{17}{64}s>s/4.
\]
This exceeds twice the band half-width and the $O(r)$ tracking error.
Hence every wall attachment in both transitions is forward. The
initial and terminal baseline pieces are also forward, since the
walls at $k+1$ and $k+T_*-1$ have disjoint ordered baseline intersection
intervals. Every edge of these candidates lies in $\mathcal T_t$.

We use incoming records indexed by their primary starting cut:
$\Gamma_j^-$ starts at $c_j$ and ends at spare cut $j+T_*$. Outgoing
records are indexed by the last primary slot they replace:
$\Gamma_j^+$ starts at spare cut $j+1-T_*$ and ends at $c_{j+1}$.
Let their normalized total corrected costs be $C_j^-$ and $C_j^+$,
using $g_j$ or $g_{j+1}$ at the primary end and $H_s$ at the spare end.
In calculating them, insert the primary compatible gauge $H_p$,
the joint world's baseline and wall gauges, and the spare compatible
gauge $H_s$ at the seams. All intermediate potentials telescope.
Every gauge change is included in $C_j^\pm$.

Here is the seam estimate at a random primary intersection. Let
$g_{\rm base}^{\rm p}$ and $g_{\rm base}^{\rm t}$ denote the gauges
of the common primary baseline in the primary and transition worlds.
Let $I_i$ be all baseline occurrences with nominal parameter in
$i+[-Cr,Cr]$, so $c_i\in I_i$. Then
\[
\begin{split}
 |g_i-g_{\rm base}^{\rm t}(c_i)|
 &\le |g_i-g_{\rm base}^{\rm p}(c_i)|\\
 &\quad+\sup_{c\in I_i}|g_{\rm base}^{\rm p}(c)-H_p(c)|
       +\sup_{c\in I_i}|H_p(c)-g_{\rm base}^{\rm t}(c)|.
\end{split}
\]
The switch envelope bounds the first term in $L^1$ by $C\beta$;
Lemma~\ref{found:lem-common-gauge} bounds the other two. These are
bounds under the full coupling, obtained from the respective world
marginals and uniform over the occurrence window. They require no
conditioning on the selected $c_i$. The other seams are treated in
the same way.

For each $\delta>0$ there is a finite $K_\delta$ such that the local
tests
\begin{equation}
 U_j^\pm=\mathbf1_{\{|C_j^\pm|\leq K_\delta\}}
 \quad\hbox{satisfy}\quad
 \mathbb P(U_j^\pm=0)\leq\delta.
 \label{improvement:transitions:eq3}
\end{equation}
To prove this, first obtain a uniform upper-tail bound for each link.
For any fixed $K$, asymmetric clipping and the negative first moment in
(L) give
\[
 K\,\mathbb P(X>K\beta)
 \leq \mathbb Ef_{J,K}(X/\beta)
       +\mathbb E(X/\beta)^-
 \leq \eta+C_X.
\]
There are $T_*-2$ links, so a union bound makes all their normalized
costs at most a sufficiently large fixed $K$ with probability at least
$1-\delta/4$. The remaining forward baseline portions, interpolated
wall movements, and all the gauge seams are bounded in absolute value
by a fixed sum of local variables having mean $C\beta$. Markov's
inequality bounds this sum by a finite multiple of $\beta$ with
probability at least $1-\delta/4$. This controls the upper tail of
$C_j^\pm$. Its negative part has bounded mean: in the joint world the
candidate is a physical path, hence its defect is nonnegative, and
(G) controls its two endpoint corrections, including the comparison
to $g_j$ or $g_{j+1}$. Markov's inequality controls the lower tail.
Increasing the numerical cap gives \eqref{improvement:transitions:eq3}.
The support and the geometry of a test do not increase when its
numerical cost cap increases.

Let $V_i$ be the absolute normalized corrected cost of the spare
baseline from its marked cut $i$ to cut $i+1$, using $H_s$. Then
\begin{equation}
 V_i\geq0,\qquad \mathbb EV_i\leq C_V,
 \qquad (V_i)_{i\in\mathbb Z}\ \hbox{is independent of}
                  \ (Z_i)_{i\in\mathbb Z}.
 \label{improvement:transitions:eq4}
\end{equation}
The transition tests may depend on both fields. Later we will bound a
successful replacement by the presence of a nearby primary flag, so the
spare-cost estimate will use the independence in
\eqref{improvement:transitions:eq4}.

Choose once and for all an integer $R\geq T_*$ large enough that,
whenever $|i-j|>R$, $Z_i$ is independent separately of $Z_j$, $U_j^-$,
and $U_j^+$. Such an $R$ follows from the explicit iid input supports.
It is independent of $s$, $\delta$, $K_\delta$, and the clipping caps.
The whole joint collection is stationary under translation by $A$.

\subsection{A finite-range replacement lemma}
\label{improvement:flags}

The remaining probabilistic step uses only the stationary fields just
constructed and their stated independence. We isolate it to make clear
how a first-moment bound suffices for the rare positive tail.

Write $F_i=\mathbf1_{\{Z_i>M\}}$, $p=\mathbb EF_0$, and
$W_i=\operatorname{clip}_M(Z_i)$. Fix
\begin{equation}
 B_*=R+1,\quad G_*=4B_*,\quad L_*=8B_*,\quad H_*=10B_*+1.
 \label{improvement:flags:eq1}
\end{equation}

\begin{lemma}[Local replacement of large positive costs]
\label{improvement:replacement}
For the stationary local fields and compatible physical pieces
constructed above, assume \eqref{improvement:primary:eq2},
\eqref{improvement:transitions:eq3}--\eqref{improvement:transitions:eq4},
and the stated independence at distances greater than $R$. There is a
stationary finite-range rule selecting disjoint replacement intervals.
It concatenates primary pieces, successful transitions, and spare
pieces. Charge each retained primary cost to its own slot and each
replacement's total normalized corrected cost to its leftmost flag.
The resulting integrable local charge $Z_i^{\mathrm{out}}$ satisfies
\begin{equation}
 \mathbb E Z_0^{\mathrm{out}}
 \le \mathbb EW_0+H_*Mp
 +C_0\{2G_*p+2(L_*+1)\delta\}
 +(H_*C_V+2K_\delta)p.
 \label{improvement:flags:eq8}
\end{equation}
The range of the rule depends on $R$ and $T_*$, independently of $M$,
$\delta$, and $K_\delta$.
\end{lemma}

\begin{proof}
\emph{The local rule.}
Consecutive flagged indices belong to the same cluster when their
difference is at most $G_*$. A cluster with extremes $a\leq b$ is
accepted exactly when $b-a\leq L_*$ and
\[
 U_{a-B_*}^-=U_{b+B_*}^+=1.
\]
For such a cluster replace the primary slots shown in
Figure~\ref{improvement:fig-replacement}:
\begin{equation}
 I(a,b)=[a-B_*,b+B_*]\cap\mathbb Z.
 \label{improvement:flags:eq2}
\end{equation}

Use the incoming transition, then the spare baseline, then the outgoing
transition. The spare portion runs from parameter $a-B_*+T_*$ to
parameter $b+B_*+1-T_*$. This interval has nonnegative length because
$B_*>T_*$. The endpoints of the whole replacement are precisely
$c_{a-B_*}$ and $c_{b+B_*+1}$. Its normalized corrected cost
$C(a,b)$ is in their primary gauge convention and satisfies
\begin{equation}
 |C(a,b)|\leq2K_\delta+\sum_{i\in I(a,b)}V_i.
 \label{improvement:flags:eq3}
\end{equation}
This follows by exact telescoping through the two transitions and the
spare gauge. Different accepted intervals are disjoint because
distinct clusters have gap greater than $G_*>2B_*$. All long or infinite
clusters are left unchanged.

\begin{figure}[htbp]
\centering
\begin{tikzpicture}[x=1cm,y=1cm,font=\small,
  route/.style={draw=MidnightBlue,line width=1pt,
    -{Stealth[length=2mm]}}]
 \draw[black!35,dashed,line width=0.8pt] (0,0)--(10,0);
 \draw[route] (-0.8,0)--(0,0);
 \draw[route] (10,0)--(10.8,0);
 \draw[route] (0,0) .. controls (0.7,0.1) and (1.1,1.6) .. (2,1.6);
 \draw[route] (2,1.6)--(8,1.6);
 \draw[route] (8,1.6) .. controls (8.9,1.6) and (9.3,0.1) .. (10,0);
 \fill[MidnightBlue] (0,0) circle (1.6pt) (2,1.6) circle (1.6pt)
   (8,1.6) circle (1.6pt) (10,0) circle (1.6pt);
 \fill[BrickRed] (3,0) circle (1.8pt) (7,0) circle (1.8pt);
 \node[below] at (3,0) {$a:\ Z_a>M$};
 \node[below] at (7,0) {$b:\ Z_b>M$};
 \node[above] at (2,1.7) {$a-B_*+T_*$};
 \node[above] at (8,1.7) {$b+B_*+1-T_*$};
 \node[above] at (5,1.6) {spare interval};
 \node[align=center] at (0.25,1.15) {incoming\\$T_*$ slots};
 \node[align=center] at (9.8,1.15) {outgoing\\$T_*$ slots};
 \node[below] at (0,-0.15) {$c_{a-B_*}$};
 \node[below] at (10,-0.15) {$c_{b+B_*+1}$};
 \draw[black!50] (0,-0.65)--(0,-0.85)--(10,-0.85)--(10,-0.65);
 \node[below] at (5,-0.85) {replaced primary slots $I(a,b)$};
\end{tikzpicture}
\caption{An accepted replacement, shown in nominal-parameter
coordinates. The marked primary slots are the extreme flags $a,b$ of
the accepted cluster; they may coincide. The flags use the matched cap
$M$, while the arrows show the actual uncapped physical route. The
diagram is schematic and does not represent a sampled FPP path; each
transition and the intervening spare piece is traversed forward.}
\label{improvement:fig-replacement}
\end{figure}

The latter instruction has an exact finite implementation. A proposed
leftmost flag $a$ and rightmost flag $b\in[a,a+L_*]$ form a complete
short cluster if and only if they are flagged, all consecutive flag
gaps in $[a,b]$ are at most $G_*$, and there is no flag in
$[a-G_*,a-1]$ or $[b+1,b+G_*]$. These tests read only
$[a-G_*,a+L_*+G_*]$. They can never accept a portion of a longer or
infinite cluster. Let $O_i$ indicate replacement of slot $i$. A cluster
covering $i$ has
\[
 a\in[i-L_*-B_*,i+B_*],
\]
so $O_i$ is determined by flags within distance $13B_*$ of $i$ and
transition tests within distance $10B_*$ of $i$. It follows in particular
that
\begin{equation}
 O_i\leq\sum_{a=i-L_*-B_*}^{i+B_*}F_a,
 \qquad \mathbb EO_i\leq H_*p.
 \label{improvement:flags:eq4}
\end{equation}

\emph{Uncovered positive costs.}
If a flagged slot is left uncovered because
its cluster is long or infinite, it contains a flag at distance
strictly greater than $R$. Following its chain from $i$ until the first
such flag gives an index $j$ with
\[
 R<|j-i|\leq R+G_*.
\]
There are $2G_*$ possible witnesses. If the cluster is short but not
accepted, one of the possible incoming tests indexed in
$[i-L_*-B_*,i-B_*]$, or outgoing tests indexed in
$[i+B_*,i+L_*+B_*]$, fails. Each of these at most $2(L_*+1)$ tests is
independent of $Z_i$, because its index is more than $R$ away. Applying
independence to these deterministic possible witnesses, before any
random choice of an extreme, gives
\begin{equation}
 \mathbb E[Z_i^+F_i(1-O_i)]
 \leq C_0\{2G_*p+2(L_*+1)\delta\}.
 \label{improvement:flags:eq5}
\end{equation}
For example each remote-flag term is
$\mathbb E[Z_i^+F_iF_j]=\mathbb E[Z_i^+F_i]p\leq C_0p$;
the failed-test terms are analogous. Thus the expectation of an
uncovered large cost is bounded directly, even if rare values carry a
substantial part of the first moment.

\emph{Retained and discarded primary costs.}
The retained costs satisfy
\begin{equation}
 (1-O_i)Z_i\leq W_i+MO_i+Z_i^+F_i(1-O_i).
 \label{improvement:flags:eq6}
\end{equation}
Indeed on a covered slot the right side is at least $-M+M=0$;
on an uncovered slot only the portion of $Z_i$ above $M$ can increase
its clipped value. In particular discarding negative primary costs
loses at most $M$ times the covered density, whether or not those
costs are correlated with nearby flags.

\emph{Replacement costs.}
By \eqref{improvement:flags:eq4} and independence of the entire spare
and primary fields,
\begin{equation}
 \mathbb E[V_iO_i]\leq H_*C_Vp.
 \label{improvement:flags:eq7}
\end{equation}
We have bounded the success indicator by the event of a nearby primary
flag; success itself need not be independent of $V_i$.

Accepted clusters have density at most $p$, so their two successful
transitions contribute at most $2K_\delta p$ to the mean upper bound.
Charge the remaining bound in \eqref{improvement:flags:eq3} to the
leftmost flag of its cluster. By stationarity, its mean per slot is
bounded by $\mathbb E[V_0O_0]$; unused spare slots in that display are
simply overcharged. Summing these bounds with
\eqref{improvement:flags:eq5}--\eqref{improvement:flags:eq7} proves
\eqref{improvement:flags:eq8}. Each charge and each reassignment involves
a uniformly bounded number of slots. Hence all charges are integrable,
and the expectation calculation requires only first moments.
\end{proof}

\subsection{Physical parameterization and endpoint accounting}
\label{improvement:locality}

The replacement lemma bounds stationary charges assigned to whole
pieces. Three further facts are needed to obtain an admissible road:
the physical concatenation must track the nominal line, its actual
period cuts must have the same mean excess as these charges, and its
full input support must remain bounded. We verify them in this order.

\paragraph{Tracking.}
Each crossing has horizontal drawdown at most $s/20$. The wall and
baseline attachments add only $O(r)$ to this drawdown; successive
walls have horizontal separation $a_0$. On each unchanged primary
slot, or each transition, take the running maximum of the horizontal
coordinate and apply the affine adjustment sending its endpoint
values to the prescribed endpoint parameters. This is continuous and
nondecreasing; constant portions are allowed. The adjustment moves
the horizontal nominal position by $O(s)$, since the horizontal
endpoint errors and the drawdown are $O(s)$. On spare portions use
the original spare nominal parameter. These parameterizations agree
at every concatenation and cover the real line.

Throughout a primary slot the vertical displacement from the nominal
line is $O(s)$. Throughout a transition it lies between its two baseline
offsets, up to $O(s)$, by
\eqref{improvement:transitions:eq1}--\eqref{improvement:transitions:eq2}.
Spare pieces have offset $10s$. Thus in the $(v,w)$ tracking norm the
entire output has tracking error at most $C_{\rm tr}s$, with
$C_{\rm tr}$ independent of the cluster and probability parameters.
Local edge repetitions and nominal pauses retain their occurrence
labels; no negative signed traversal is ever substituted for physical
motion.
Every baseline--wall and wall--link splice is at a lattice vertex.
A marked spare cut inside an edge only subdivides consecutive forward
travel along that same baseline occurrence. Merging those subdivisions
therefore leaves a two-sided lattice walk with locally finite edge
occurrences, as required by Definition~\ref{def:fb-road}.

\paragraph{Actual period cuts and integrability.}
To convert the charge bound \eqref{improvement:flags:eq8} into a mean
actual excess, assign the primary gauges to the designated primary
concatenation endpoints, and the compatible spare gauges to spare
endpoints. On each primary or transition piece interpolate its two
endpoint gauges linearly in the fraction of its positive physical
cost already traversed; on spare pieces use the spare interpolation.
These assignments agree at concatenations, are local, and commute
with translation by $A$. In particular they assign a gauge at every
actual first occurrence of an integer nominal level, including when
that occurrence precedes the designated endpoint during a nominal
pause. The assigned gauges are integrable at all these cuts.
Integrability on a transition follows from its positive physical
costs. On success, the total physical passage cost is at most
\[
 K_\delta\beta+|g_{\rm first}|+|g_{\rm last}|+C n,
\]
where the appropriate endpoint gauge may be $H_s$. The right side is
integrable by (G), and any physical subpiece costs at most the whole
physical path. A selected primary candidate has physical cost no
larger than its baseline competitor: their endpoints and corrections
are identical. That competitor lies in a fixed number of integrable
input periods. On primary and transition pieces, linear interpolation
bounds each assigned gauge by the absolute endpoint gauges. On spare
pieces, the original interpolation and (G) give integrability over
bounded parameter windows. Partial primary pieces, successful
transitions, and spare periods are therefore integrable.
This integrability is required at each fixed $n$; the bound need not be
uniform in $n$.

For an interval of $N$ slots, summing the local upper charges used in
\eqref{improvement:flags:eq8} bounds the corrected cost of all complete
primary and replacement pieces. Passing to the actual first-occurrence
endpoint cuts changes this bound only within the bounded number of
primary pieces and replacement intervals meeting the two boundaries.
The primary boundary pieces must be included even when no replacement
is accepted: a terminal wall arm may follow the first occurrence of
the terminal integer. Each boundary remainder is bounded by a fixed
finite sum of the integrable costs and gauges just described. Its
expected absolute value is bounded independently of $N$.
After division by $N$ this boundary error vanishes. The endpoint gauge
difference has zero mean: the complete output rule, including the
choice of endpoint label, is stationary under $A$, so the gauges at
translated cuts have the same integrable distribution. Thus
\eqref{improvement:flags:eq8} indeed bounds the expected actual $Q_n$
per old slot, divided by $\beta$. This uses stationarity of the chosen
labels and does not assume that the mean of the calibration at an
arbitrary random point is zero.

\paragraph{Input locality.}
The finite cluster test and the input supports in
Section~\ref{improvement:worlds} imply a fixed output dependence
radius $R_{\rm out}$ in old spatial units. Numerical caps affect only
comparisons of already calculated local quantities. The large
auxiliary boxes used to prove (L), and the size of its finite cap set,
are not inputs to this path construction. Neither $R_{\rm out}$ nor
$C_{\rm tr}$ depends on $s$, the directional tolerances, the cap, the
seam failure probability, or the input road rule.

\subsection{Choice of constants and contraction of the infimum}
\label{improvement:contraction}

The construction can now be closed in the class used to define $b(n)$.
The order of choices is essential: choose the rescaling from structural
support and tracking bounds first, then the small inward tilt, and only
then the probability and cap tolerances.

\begin{proposition}[Contraction in the original road class]
\label{improvement:contraction-proposition}
For suitable fixed rational $v,w$ near the flat endpoints and $s>0$,
there are an integer $D\ge1$ and $\rho<1$ such that
\[
 b(Dn)\le\rho b(n)
\]
for every sufficiently large admissible $n$. The infimum on both sides
uses the road class of Definition~\ref{def:fb-road}, with costs in their
original units.
\end{proposition}

\begin{proof}
The structural constants $C_{\rm tr},R_{\rm out},C_1$ are uniform
in sufficiently small $s$ and the nearby rational directions.
First choose an integer
\begin{equation}
 D\geq\max\{1000C_{\rm tr},R_{\rm out}/10,1\}.
 \label{improvement:contraction:eq2}
\end{equation}
Next fix rational $0<s<1/10$ and $\epsilon>0$ so small that
$DC_1s<1/4$ and $D\epsilon<1/4$. The geometry and the constants
$R,B_*,G_*,L_*,H_*,C_0,C_V$ have already been fixed uniformly.
If $C_0=0$, the primary costs vanish almost surely and no replacements
are necessary. Otherwise choose
\[
 \delta\leq\min\left\{\frac12,
       \frac{\epsilon}{6(L_*+1)C_0}\right\},
\]
then the corresponding transition cap $K_\delta$. Set
\[
 A_*=2G_*C_0+H_*C_V+2K_\delta,
 \qquad
 M_0\geq\max\{1,3C_0A_*/\epsilon\},
 \qquad \eta\leq\frac{\epsilon}{3(H_*+1)}.
\]
Choose the finite matched-cap range from (L) with this $M_0,\eta$,
and then fix the geometric comparison tolerances and the rational
directions so that it applies uniformly over that range. At every
sufficiently large admissible $n$, choose its deterministic good cap.
Equation~\eqref{improvement:primary:eq4} gives $Mp\leq\eta$, while
$p\leq C_0/M_0$ by Markov's inequality. In
\eqref{improvement:flags:eq8} the failed-test error, the terms
$A_*p$, and the terms $(H_*+1)\eta$ are each at most
$\epsilon/3$. We obtain an actual stationary local road with
\begin{equation}
 \mathbb E Q_n(\hbox{one old output slot})
       \leq(C_1s+\epsilon)\beta.
 \label{improvement:contraction:eq1}
\end{equation}
Regard $D$ old slots as one new period and use space units $Dn$.
The tracking error is then at most $s/1000$, the input radius at most
$10$, and the rule commutes with translation by $DA$ in old units,
that is by $A$ in new units. Reindexing the private product field by
this rescaling preserves its shift-invariant law: a new lattice shift
corresponds to the old lattice shift multiplied by $D$. Thus it is in the original road class at
scale $Dn$. The same construction in the other direction gives a pair
with common mean bound
\begin{equation}
 D(C_1s+\epsilon)\beta.
 \label{improvement:contraction:eq3}
\end{equation}
All the rational meshes have fixed denominators; multiplication of
an admissible $n$ by the integer $D$ preserves admissibility.

Let $b(n)$ be the infimum of the common mean bound over that road
class. Invoke the uniform input estimates with the fixed enlarged
prefactor $2C_s$; this changes only the threshold for $n$. That
threshold is therefore uniform over members with
$\beta\leq2C_s n^{1-c}$. Since $b(n)\leq C_s n^{1-c}$, for each
sufficiently small $t>0$ one can choose a pair with
$\beta\leq b(n)+t\leq2C_s n^{1-c}$. Apply
\eqref{improvement:contraction:eq3} and then let $t\downarrow0$; existence
of a minimizing road is unnecessary. This proves
\begin{equation}
 b(Dn)\leq D(C_1s+\epsilon)b(n).
 \label{improvement:contraction:eq4}
\end{equation}
All directions, meshes, and tolerances are fixed in the stated order,
and the choices give $\rho=D(C_1s+\epsilon)<1/2$. This proves the
assertion for all admissible $n$ above the common threshold.
\end{proof}
